% Planner section, 2026-11-06. Sensitivity checks and scope notes prompted by two AI reviews.
% Numbers: docs/track_r/planner_analysis/A26_reviewer_checks/out/REVIEWER_CHECKS_20261106.json (ledger C-V1 to C-V3).
\section{Sensitivity checks and scope}
\label{sec:suppl-sensitivity}

This section answers questions a reader may ask about our choices.
The checks are exploratory and post hoc.
They reuse stored runs; no new run was drawn.

\subsection{The \good threshold}
\label{sec:suppl-theta}
The thresholds $\theta$ ($25$\,dB on FFHQ, $20$\,dB on ImageNet) sit below the PSNR of typical correct runs and above most failures.
For SITCOM on FFHQ, the \good runs have a median PSNR of $27.3$\,dB, with $10$th and $90$th percentiles of $26.5$ and $28.0$\,dB; its \cons failures have a median of $15.4$\,dB.
Only $3.5\%$ of its $2{,}880$ runs lie within $1$\,dB of $\theta$.
\Cref{tab:suppl-theta-sitcom,tab:suppl-theta-pcdaps} rescore the policies at other thresholds, on the panels whose per-run PSNR is stored.

\begin{table}[h]
\centering\small
\setlength{\tabcolsep}{4pt}
\begin{tabular}{@{}c cccc c ccc@{}}
\toprule
 & \multicolumn{4}{c}{$K=4$} & & \multicolumn{3}{c}{$K=16$} \\
\cmidrule(lr){2-5}\cmidrule(l){7-9}
$\theta$ (dB) & First & $\Eo$ & $\Eone$ & $\Oorb$ & & $\Eo$ & $\Eone$ & $\Oorb$ \\
\midrule
22 & 63.6 & 83.8 & 90.4--90.8 & 91.1 & & 86.7 & 96.1 & 96.7 \\
23 & 62.6 & 82.7 & 89.2--89.6 & 89.9 & & 85.6 & 95.0 & 96.1 \\
24 & 60.5 & 80.5 & 86.5--86.9 & 87.0 & & 83.9 & 92.8 & 93.9 \\
\textbf{25} & 58.6 & 78.4 & 84.2--84.5 & 84.9 & & 82.2 & 91.1 & 92.2 \\
26 & 56.9 & 75.6 & 80.7--81.0 & 82.0 & & 78.3 & 86.1 & 88.9 \\
27 & 39.4 & 49.0 & 52.3--52.6 & 54.1 & & 50.0 & 55.0 & 57.8 \\
\bottomrule
\end{tabular}
\caption{\textbf{SITCOM on FFHQ ($180$ images, $16$ runs each): \% of images with a \good output at other thresholds.} Exact averages over all $K$-subsets of each image's runs. $\Eone$ uses the face test. For $34$ of the $2{,}880$ runs, the test turns the run although its best orbit member is not $R_{111}$; the PSNR of that rotation is not stored, so $\Eone$ is bracketed by counting these runs as failures or successes. Bold: the registered threshold.}
\label{tab:suppl-theta-sitcom}
\end{table}

\begin{table}[h]
\centering\small
\setlength{\tabcolsep}{3.5pt}
\begin{tabular}{@{}c cccc c c cccc@{}}
\toprule
\multicolumn{5}{c}{FFHQ ($180$ images)} & & \multicolumn{5}{c}{ImageNet ($100$ images)} \\
\cmidrule(r){1-5}\cmidrule(l){7-11}
$\theta$ & Random & $\Eo$ & $\Eone$ & $\Oraw$ & & $\theta$ & Random & $\Eo$ & $\Eone$ & $\Oorb$ \\
\midrule
22 & 83.6 & 86.7 & 88.9 & 87.8 & & 17 & 40.2 & 43.0 & 63.0 & 69.0 \\
23 & 81.4 & 85.6 & 87.8 & 86.1 & & 18 & 38.2 & 42.0 & 62.0 & 67.0 \\
24 & 80.3 & 85.0 & 87.2 & 85.6 & & 19 & 33.2 & 39.0 & 58.0 & 63.0 \\
\textbf{25} & 79.0 & 82.8 & 84.4 & 83.3 & & \textbf{20} & 30.8 & 37.0 & 53.0 & 55.0 \\
26 & 77.8 & 81.7 & 83.3 & 82.2 & & 21 & 29.2 & 34.0 & 49.0 & 51.0 \\
27 & 76.1 & 78.9 & 80.6 & 81.1 & & 22 & 25.8 & 30.0 & 46.0 & 48.0 \\
28 & 74.9 & 78.9 & 80.6 & 79.4 & & 23 & 25.2 & 29.0 & 44.0 & 45.0 \\
\bottomrule
\end{tabular}
\caption{\textbf{PC-DAPS, one sample of four particles per measurement: \% of images with a \good output at other thresholds.} Random: the mean over particles, which PC-DAPS reports. $\Eone$ uses the face test on FFHQ and the generic test on ImageNet. The orbit oracle is not stored for the FFHQ particles; $\Oraw$ is shown instead. Bold: the registered thresholds.}
\label{tab:suppl-theta-pcdaps}
\end{table}

The comparisons of the main text hold across these thresholds\claim{C-V1}.
For SITCOM, $\Eone$ stays within $1.8$ points of the orbit oracle at $K=4$ and within $2.8$ points at $K=16$, from $22$ to $27$\,dB.
All policies drop at $27$\,dB, because SITCOM's correct runs cluster just above it.
Choosing among PC-DAPS particles adds $4.5$ to $6.9$ points over a random particle on FFHQ, from $22$ to $28$\,dB, and $18.8$ to $24.8$ points on ImageNet, from $17$ to $23$\,dB.
On ImageNet, $\Eone$ stays within $6$ points of the orbit oracle at every threshold.
The mean PSNR, SSIM and LPIPS results of the main text do not use $\theta$ at all.

\subsection{The label-free threshold and the pilot size}
\label{sec:suppl-pilot}
The label-free rule takes the per-image minimum residual over a pilot of runs and returns the median plus three scaled median absolute deviations.
\Cref{tab:suppl-pilot} compares it with the labeled threshold for pilots of two, three and four runs (the first runs of each image).
At $\sigma=0.05$, it lands within $2.4\%$ even with two runs\claim{C-V2}.
The largest miss is on the enriched cohort~B at $\sigma=0.02$ ($+7.6\%$ with two runs, $+3.5\%$ with four), where a few hard images make the minima heavy-tailed.
Medians resist such tails, and they need no labels, so the rule can also be updated as images arrive.

\begin{table}[h]
\centering\small
\begin{tabular}{@{}lcccc@{}}
\toprule
Solver, data, cohort & $\sigma$ & 2 runs & 3 runs & 4 runs \\
\midrule
SITCOM, FFHQ, A & $0.05$ & +1.7 & +1.7 & +1.8 \\
SITCOM, FFHQ, B & $0.05$ & +2.4 & +2.2 & +1.9 \\
DAPS-4K, FFHQ, A & $0.05$ & +0.3 & +0.2 & +0.2 \\
DAPS-4K, FFHQ, B & $0.05$ & +0.6 & +0.4 & +0.3 \\
DAPS-4K, ImageNet & $0.05$ & +1.4 & +1.2 & +1.0 \\
FAST-DIPS, FFHQ, A & $0.02$ & +0.5 & +0.1 & +0.1 \\
FAST-DIPS, FFHQ, B & $0.02$ & +7.6 & +4.9 & +3.5 \\
FAST-DIPS, FFHQ, A & $0.10$ & +0.6 & +0.7 & +0.6 \\
FAST-DIPS, FFHQ, B & $0.10$ & +1.1 & +1.0 & +1.2 \\
\bottomrule
\end{tabular}
\caption{\textbf{Label-free threshold against the labeled one} ($95$th percentile of $\rho$ over \good runs), in \% of the labeled value, by pilot size.}
\label{tab:suppl-pilot}
\end{table}

\subsection{Dependence between runs}
\label{sec:suppl-dependence}
Each run uses a fresh seed and shares only the measurement with the other runs; no run is warm-started or tuned on another.
The runs of an image are therefore exchangeable.
By de Finetti's theorem, exchangeable runs behave as independent draws given the image's class rates.
This is what \cref{M-prop:stop} and the Dirichlet mixture of \cref{M-sec:selection} assume: the image's difficulty is the only link between its runs.
Runs of an easy image often agree closely, because its posterior is concentrated; this reflects the image, not the seeds.
Exchangeability would fail if outcomes depended on the run index.
We found no drift across run blocks (\cref{M-sec:setup}).
Adjacent runs also agree no more often than any two runs of the same image\claim{C-V3}.
On five panels (SITCOM, FFHQ, $16$ runs; DAPS-4K, FFHQ and ImageNet, $8$ runs; FAST-DIPS, FFHQ, $\sigma=0.02$ and $0.10$, $8$ runs), the concordance of the \good indicator for adjacent runs minus that for all pairs lies between $-0.010$ and $+0.012$.
The permutation $p$-values, with the run order shuffled within images, are $0.17$ to $0.69$.
PC-DAPS particles interact, so a PC-DAPS sample is not four independent runs; we never apply the stopping rule within it.

\subsection{Compute}
\label{sec:suppl-compute}
Within a solver, every run makes the same number of denoiser calls: $4{,}000$ for DAPS-4K and $2{,}000$ for DAPS-2K, for instance.
The equal-mean-runs comparisons of \cref{M-sec:protocol} are therefore equal-compute comparisons.
The protocol adds one evaluation of the operator per run, an FFT of the $384\times384$ canvas per channel, and one orientation test per output ($8.5$\,ms with CLIP).
On one otherwise idle GPU, a DAPS-4K run takes $66.0$\,s and a PC-DAPS sample $30.6$\,s (\cref{tab:s4-pcdaps}).
We compare solvers at equal compute only for PC-DAPS against DAPS-4K, where calls and seconds are reported.

\subsection{Other verifiers}
\label{sec:suppl-verifiers}
Under Gaussian noise, the likelihood of a run depends on $y$ only through its residual $\rho$.
A reweighted or robust residual is the likelihood of another noise model.
It can help when the noise model is wrong, for instance with outliers.
It cannot separate a \cons run from a \good one in principle, because both fit the measurement within the noise.
Such runs differ in the image domain, where only the prior can rank them.

\subsection{Orientation scorers}
\label{sec:suppl-orientation}
Within an orbit, the Bayes choice orders the members by the prior density (\cref{M-sec:setup}).
CLIP prompts, the paired denoising-loss score and a learned upright-versus-rotated classifier are three proxies for this order.
A learned classifier needs training images from the target domain; we did not train one, to keep the protocol zero-shot.
Any proxy plugs into step~(3) of the protocol unchanged.
The gap decomposition of \cref{M-sec:audit} bounds what a better proxy can gain: the orientation part of the gap to the orbit oracle, $4.0$ points for DAPS-4K at $K=8$ and $5.0$ points for FAST-DIPS at $K=64$ on ImageNet, and at most $0.6$ points on FFHQ.

\subsection{Other operators, tasks and solvers}
\label{sec:suppl-scope}
The residual comparison of \cref{M-sec:protocol} needs only a fixed measurement with Gaussian noise.
\cref{M-prop:stop} needs only exchangeable runs and a pass event.
Neither depends on the operator or the solver.
What depends on them is the mix of failures, and so which component helps.
We state our expectations; we did not test them.
\begin{itemize}[leftmargin=*,itemsep=1pt,topsep=2pt]
\item \emph{Coded diffraction} removes the twin ambiguity. Orbit errors disappear, and step~(3) becomes unnecessary.
\item \emph{Lower oversampling} admits more images that fit the data, so we expect more \cons failures and less benefit from stopping. Higher oversampling should act the other way.
\item \emph{Inpainting and super-resolution} have large null spaces. A wrong run can then fit the data exactly. We expect most failed runs to be \cons, so stopping should not help and selection must rely on the prior.
\item \emph{Other solvers}, including flow-based ones, enter through their runs only. The protocol needs a residual and a threshold, both available for any solver.
\item \emph{Real measurements} have no ground truth. The protocol needs none, but scoring it does; we therefore use simulated measurements.
\end{itemize}

\subsection{Forecast scores over all six tests}
\label{sec:suppl-scores}
Coverage counts reward wide intervals.
\Cref{tab:suppl-scores} therefore also gives, for every registered forecast of the selected run, the absolute error of the point forecast, the interval width and the $95\%$ interval score~\cite{gneiting2007scoring}: the width plus $40$ times the distance by which the outcome misses the interval.
These are end-to-end prospective forecasts: each was fitted to the runs observed before its freeze and scored on later runs\claim{C-R8}.
The conditional analysis of \cref{M-sec:selection}, which uses each image's realized classes, is a retrospective diagnosis and is not scored here.

\begin{table}[h]
\centering\small
\setlength{\tabcolsep}{4pt}
\begin{tabular}{@{}cl cccc c cccc@{}}
\toprule
 & & \multicolumn{4}{c}{Race} & & \multicolumn{4}{c}{(S2)} \\
\cmidrule(lr){3-6}\cmidrule(l){8-11}
Test & Solver, data & Cov. & Error & Width & IS & & Cov. & Error & Width & IS \\
\midrule
1 & FAST-DIPS, FFHQ & 9/42 & 4.4 & 5.9 & 60.9 & & 36/42 & 1.8 & 8.4 & 11.4 \\
2 & FAST-DIPS, ImageNet & 9/9 & 4.2 & 16.2 & 16.2 & & 0/9 & 16.5 & 13.7 & 397.7 \\
3 & SITCOM, FFHQ & 18/18 & 1.7 & 12.4 & 12.4 & & 2/18 & 14.3 & 12.6 & 338.6 \\
4 & DAPS-2K, FFHQ & 12/14 & 2.1 & 6.7 & 9.5 & & 14/14 & 1.7 & 7.8 & 7.8 \\
5 & FAST-DIPS, ImageNet & 18/18 & 1.6 & 10.7 & 10.7 & & 1/18 & 11.6 & 11.0 & 263.8 \\
6 & FAST-DIPS, ImageNet$'$ & 10/10 & 3.5 & 11.6 & 11.6 & & 0/10 & 17.3 & 9.3 & 516.4 \\
\midrule
All & & 76/111 & 3.1 & 9.2 & 30.3 & & 53/111 & 8.0 & 10.0 & 181.8 \\
\bottomrule
\end{tabular}
\caption{\textbf{Registered forecasts of the selected run ($\Eo$, $\Eone$): coverage, mean absolute error, mean width and mean interval score, in points.} Means are over forecasts, so tests with more $K$ values weigh more. For the oracles, the two models share point forecasts (mean error $3.8$ points, biased upward by $3.2$); their interval scores are $48$ (race) and $45$ (S2).}
\label{tab:suppl-scores}
\end{table}

The race loses test~1 clearly and test~4 narrowly.
On test~1 its selected-run forecasts are $4.4$ points too high on average, because the shared generation component expected too many images to produce a \good run (\cref{M-sec:selection}).
A forecaster with a separate near-zero-success component may do better there; we did not register one.

\subsection{Pass probability and class probability}
\label{sec:suppl-kappa}
\Cref{M-prop:stop} uses the probability $s$ that a run passes $\rho\le\tau_S$.
The forecasts use $\kappa$, the probability that a run is not \incons.
They differ in two ways\claim{C-V5}.
First, by construction about $5\%$ of \good runs exceed the labeled threshold: $82$ of SITCOM's $1{,}688$ \good runs on FFHQ, and $11$ of its $94$ orbit errors.
Second, a label-free $\tau_S$ that lies slightly above the labeled $\tau$ lets some \incons runs pass: $40$ of $350$ for DAPS-4K on FFHQ, and $139$ of $210$ on ImageNet, where the residuals of failed runs sit just above $\tau$.
Labels are therefore tied to the labeled $\tau$, while stopping uses $\tau_S$.

\subsection{How accurately must $\tau_S$ be set?}
\label{sec:suppl-tauacc}
Selection ($\Eo$, $\Eone$) does not use $\tau_S$; only stopping and $\textsc{flag}_C$ do.
\Cref{tab:suppl-tauacc} reruns the frozen protocol on the held-out panels with other thresholds and other targets $M$\claim{C-V4}.
For DAPS, residuals are concentrated so tightly that a threshold $2\%$ too low lets no run pass.
The rule then never stops: it spends $N_{\max}$ runs and flags every image, which a user sees without labels.
A threshold $2\%$ too high lets almost every run pass, so the rule stops after $M$ runs and keeps only part of the gain.
The labeled threshold, the label-free threshold and that of the other cohort's pilot lie within $0.4\%$ of each other on these panels, and all keep the gain ($+2.6$ to $+5.0$ points on FFHQ).
On the DAPS-4K panels, the label-free pilot reused the first four evaluated runs of the same images.
A user would instead run it once per solver and noise level; its cost is then amortized and is not part of the mean runs per image.
The cross-cohort column shows that a pilot on other images works as well.

\begin{table}[h]
\centering\small
\setlength{\tabcolsep}{3.5pt}
\begin{tabular}{@{}l cc cc ccc@{}}
\toprule
 & \multicolumn{2}{c}{Calibrated $\tau_S$} & \multicolumn{2}{c}{Registered $\tau_S$ scaled} & \multicolumn{3}{c}{Target $M$} \\
\cmidrule(lr){2-3}\cmidrule(lr){4-5}\cmidrule(l){6-8}
Solver, data, cohort & labeled & label-free & $\times0.98$ & $\times1.02$ & 1 & 3 & 4 \\
\midrule
DAPS-2K, FFHQ, A & $+3.0$ & $+3.3$ & $\phantom{+}0.0$ (8.0) & $+1.3$ (2.1) & $+8.4$ & $+1.4$ & $+0.8$ \\
DAPS-2K, FFHQ, B & $+3.9$ & $+4.0$ & $\phantom{+}0.0$ (8.0) & $+0.7$ (2.5) & $+7.2$ & $+2.1$ & $+1.2$ \\
DAPS-4K, FFHQ, A & $+2.6$ & $+3.1$ & $\phantom{+}0.0$ (8.0) & $+1.2$ (2.2) & $+11.2$ & $+1.1$ & $+0.5$ \\
DAPS-4K, FFHQ, B & $+4.2$ & $+5.0$ & $\phantom{+}0.0$ (8.0) & $+1.9$ (2.5) & $+9.0$ & $+3.0$ & $+1.9$ \\
DAPS-4K, ImageNet & $+1.3$ & $+0.1$ & $+2.2$ (6.4) & $+0.2$ (2.0) & $+1.0$ & $-0.1$ & $-0.1$ \\
SITCOM, FFHQ, A & $-0.1$ & $+1.6$ & $-1.5$ (11.5) & $+1.5$ (2.2) & $-0.3$ & $+0.0$ & $+0.1$ \\
SITCOM, FFHQ, B & $-0.9$ & $+1.0$ & $-3.8$ (9.2) & $+1.1$ (2.6) & $-2.4$ & $-0.2$ & $+0.1$ \\
\bottomrule
\end{tabular}
\caption{\textbf{Gain of the frozen protocol over fixed best-of-$K$ at equal mean runs, in points, under other thresholds and targets} (exploratory). Registered: labeled $\tau_S$ for DAPS-2K and SITCOM, label-free for DAPS-4K; $M=2$ except in the last three columns, which use the registered $\tau_S$. In parentheses: mean runs per image. $N_{\max}=8$ (DAPS) and $16$ (SITCOM). With the other cohort's label-free threshold, DAPS-4K gains $+3.0$ (A) and $+4.5$ (B).}
\label{tab:suppl-tauacc}
\end{table}

\subsection{The target $M$}
\label{sec:suppl-M}
\Cref{M-fig:teaser}b shows the protocol for $M=1$ to $4$ on the development panel.
Each $M$ is a point on a frontier that lies above fixed best-of-$K$; a larger $M$ buys accuracy with runs.
On the held-out DAPS panels (\cref{tab:suppl-tauacc}), the gain at equal mean runs is largest for $M=1$ ($+7.2$ to $+11.2$ points at $1.3$ to $2.2$ runs) and shrinks as $M$ grows, because the fixed-$K$ curve flattens.
$M=2$ was fixed before these panels were scored; we report the others as exploratory.
Requiring a confident orientation before stopping does not help on ImageNet (\cref{tab:s4-frontier}).

\subsection{The orientation flag}
\label{sec:suppl-flago}
$\textsc{flag}_O$ marks the least confident $30\%$ of a batch.
It is a triage rule: it flags a fixed share whatever the batch, and one output's flag can change when others are added.
It is not a calibrated confidence.
A deployment that needs a fixed rule can freeze an absolute margin, as the confidence-conditioned variant of \cref{tab:s4-frontier} does with the cut of the $16$-run ImageNet panel.

\subsection{Reproducing the published DAPS protocol}
\label{sec:suppl-dapsrepro}
Our ImageNet ground truth and measurements equal those of the DAPS code for all $100$ images, and every DAPS-4K run has the published configuration ($400$ annealing steps with $10$ ODE steps; $4{,}000$ denoiser calls, from the run receipts).
On our $180$ FFHQ images, the per-metric best of four gives $29.76$\,dB, SSIM $0.801$ and VGG-LPIPS $0.156$, against $30.63$\,dB, $0.851$ and $0.139$ in the DAPS paper on its own FFHQ images\claim{C-A12}.
On ImageNet we obtain $22.28$\,dB, against $25.78$\,dB in the DAPS paper and $20.98$\,dB in an independent rerun~\cite{kocabas2026pcdaps}.
We could not explain the ImageNet gap.
Our comparisons between policies always use the same runs, so the gap does not affect them; it does affect any comparison with the published number.
